% ============================================================
% Paper deep reading: TempFlow-GRPO
% ============================================================

\section{论文精读：TempFlow-GRPO}

\begin{conceptbox}
\textbf{论文：} \href{https://arxiv.org/abs/2508.04324}{TempFlow-GRPO: When Timing Matters for GRPO in Flow Models}\\
\textbf{简称：} TempFlow-GRPO。\\
\textbf{版本：} arXiv 2508.04324v4，submitted on 2025-08-06，last revised on 2025-10-15；OpenReview 为 ICLR 2026 Poster。\\
\textbf{方向标签：} Flow Matching、GRPO、Temporal Credit Assignment、Trajectory Branching、Noise-aware Weighting。\\
\textbf{阅读日期：} 2026-04-27。\\
\textbf{前置：} 已读 DenseGRPO 和 Chunk-GRPO，当前关心“step-level reward 是否必要、chunk 是否只是平滑 credit、如何避免 A+B 式组合”。
\end{conceptbox}

\subsection{0. 一句话结论}

TempFlow-GRPO 最值得关注的点不是“又一个 Flow-GRPO variant”，而是它提出了一种构造 step-level process reward 的新方式：

\[
\begin{aligned}
&\text{不是直接给中间 latent 打分}\\
&\rightarrow\quad
\text{在某个 timestep 做 controlled branching，}\\
&\hspace{2.2em}\text{再用最终 reward 反推该 step 的贡献}
\end{aligned}
\]

它的核心动作是：

\[
\begin{aligned}
\text{ODE 到 } x_k
&\rightarrow
\text{只在 } k \text{ 做一次 SDE 探索}\\
&\rightarrow
\text{ODE 到 } x_0
\rightarrow
\text{用最终 reward 归因给 } k
\end{aligned}
\]

\begin{summarybox}
\textbf{我的判断：} 这篇比 Chunk-GRPO 更直接地回答了我们刚才的问题：\key{每个 step 的 reward 可以不靠中间 reward model，而靠 controlled intervention + final reward 构造出来。} DenseGRPO 是“从中间 latent 垂直 readout”；TempFlow-GRPO 是“在指定时刻水平分叉探索，再把最终差异归因给分叉点”。
\end{summarybox}

\subsection{1. 它要解决什么}

Flow-GRPO / Dance-GRPO 的基本问题是：

\[
\hat{A}_t^i
=
\frac{
R(x_0^i,c)-\mathrm{mean}(\{R(x_0^g,c)\}_{g=1}^G)
}{
\mathrm{std}(\{R(x_0^g,c)\}_{g=1}^G)
}
\]

这里 \(\hat{A}_t^i\) 实际上没有真正的 \(t\)-specific 信息。它虽然写成 \(\hat{A}_t^i\)，但 reward 来自最终 \(x_0\)，所以整条 trajectory 的所有 timesteps 基本共用同一个 terminal advantage。

TempFlow-GRPO 指出这有两个问题：

\begin{enumerate}[leftmargin=1.5em]
  \item \textbf{sparse terminal reward：} 只知道最终结果好坏，不知道哪个 timestep 的探索带来提升。
  \item \textbf{temporal uniformity：} 所有 timesteps 被等权优化，但生成过程里不同时间点的探索能力不同。
\end{enumerate}

作者先做了一个关键观察：只在某一个 denoising step 上引入 SDE stochasticity，其余地方都用 ODE，最后看 PickScore 的 reward std。结果显示 early steps 的 reward variance 最大，late steps 接近 0。

\begin{center}
  \includegraphics[width=0.46\linewidth]{tempflowgrpo/fig2_left_reward_variance.png}
  \hfill
  \includegraphics[width=0.46\linewidth]{tempflowgrpo/fig2_right_branching.png}

  {\small \textbf{图 1：} TempFlow-GRPO 的动机。左图：只在单个 timestep 注入 SDE 噪声后，reward variance 在 early steps 最大。右图：从确定性 ODE 轨迹上分叉出 stochastic branch，并把最终 reward 差异归因给分叉点。图源：arXiv HTML。}
\end{center}

\begin{summarybox}
\textbf{这张图很重要：} 它把“不同 step 贡献不同”从直觉变成了一个可实验测量的问题。早期 step 的随机扰动会显著改变最终 reward，说明早期确实承担结构性探索；后期 step 的扰动影响小，更像细节 refinement。
\end{summarybox}

\subsection{2. 方法一：Trajectory Branching}

\subsubsection{2.1 为什么不直接训练 process reward model}

传统 process reward 的想法是直接评价中间状态 \(x_t\)。但在 diffusion / flow 模型里，中间 latent 或 partially-denoised image 很难被 reward model 稳定理解：

\[
x_t \quad \text{semantic ambiguous}
\]

所以 TempFlow-GRPO 换了一个角度：不评价 \(x_t\) 本身，而是做一个 intervention：

\[
\text{如果只在 } t=k \text{ 这个点探索，最终 reward 会怎么变？}
\]

这就把问题从“中间状态评分”变成“受控分叉后的最终结果评分”。

\subsubsection{2.2 Branching 的定义}

先从初始噪声 \(x_T\) 出发，用 ODE 确定性走到指定 timestep \(k\)：

\[
x_T \xrightarrow{\mathrm{ODE}} x_k
\]

然后只在 \(k\) 这个点引入一次 SDE stochasticity：

\[
x_{k-1}
=
\mathrm{SDE}(x_k,\epsilon)
\]

最后从 \(x_{k-1}\) 到 \(x_0\) 全部用 ODE 确定性完成：

\[
x_0
=
\mathrm{ODE}^{k-1}(x_{k-1})
\]

所以第 \(i\) 个 branch 的最终可评分样本是：

\[
\tilde{x}_{0,k}^i
=
\mathrm{ODE}^{k-1}
\left(
\mathrm{SDE}(x_k^i,\epsilon^i)
\right)
\]

对应 reward 是：

\[
R_k^i
=
R(\tilde{x}_{0,k}^i,c)
=
R\left(
\mathrm{ODE}^{k-1}
\left(
\mathrm{SDE}(x_k^i,\epsilon^i)
\right),
c
\right)
\]

然后用同一个 \(k\) 下多个 branches 的 reward 做 group normalization：

\[
\hat{A}_k^i
=
\frac{
R_k^i-\mathrm{mean}(\{R_k^g\}_{g=1}^{G})
}{
\mathrm{std}(\{R_k^g\}_{g=1}^{G})+\epsilon
}
\]

\begin{summarybox}
\textbf{关键理解：} 这个 \(\hat{A}_k^i\) 不是 DenseGRPO 的 \(\Delta R_k=R_{k-1}-R_k\)。它衡量的是：在相同 \(x_k\) 附近，不同 stochastic branch 最终会产生怎样的 reward 差异。因此它更像 \key{branch-level counterfactual advantage}。
\end{summarybox}

\subsubsection{2.3 为什么它能做 credit localization}

因为除了 \(k\) 这个点之外，其余路径都被 ODE 固定住了。也就是说，在同一个 initial noise 和 prompt 下，如果多个 branches 最后 reward 不同，差异主要来自：

\[
\epsilon_k
\quad \text{at branching timestep } k
\]

这就是它所谓的 credit localization：

\[
\text{reward variation}
\quad \Rightarrow \quad
\text{branching action at } k
\]

这和 DenseGRPO 有本质区别：

\[
\text{DenseGRPO：}
\quad
x_k \xrightarrow{\mathrm{ODE\ readout}} \hat{x}_0
\quad
\text{给 } x_k \text{ 打分}
\]

\[
\text{TempFlow-GRPO：}
\quad
x_k \xrightarrow{\mathrm{SDE\ branch}} x_{k-1}
\xrightarrow{\mathrm{ODE}} \hat{x}_0
\quad
\text{给 branch action 打分}
\]

\begin{warnbox}
\textbf{需要注意：} 它不是证明“第 \(k\) 步独立决定最终 reward”。它只是通过控制变量，让 stochastic difference 只发生在 \(k\)。后续 ODE 仍然会放大、修正或过滤这个差异，所以它更准确的说法是：\key{把探索差异定位到 branching point}，而不是完全因果地分解整条 denoising 链。
\end{warnbox}

\begin{center}
  \includegraphics[width=0.90\linewidth]{tempflowgrpo/fig5_framework.png}

  {\small \textbf{图 2：} TempFlow-GRPO 框架图。核心是 ODE-SDE-ODE 的分叉采样：在指定 timestep 注入 stochasticity，之后用最终 reward 形成该 timestep 的 process signal。图源：arXiv HTML。}
\end{center}

\subsection{3. 方法二：Noise-Aware Policy Weighting}

Trajectory branching 解决的是：

\[
\text{每个 timestep 的 advantage 怎么来？}
\]

Noise-aware weighting 解决的是：

\[
\text{每个 timestep 的更新强度应该一样吗？}
\]

答案是不应该一样。SDE sampler 在不同 timestep 的噪声注入幅度不同：

\[
\sigma_t\sqrt{\Delta t}
\]

早期噪声大，探索空间大；后期噪声小，主要是细节 refinement。作者观察到 reward std 和 noise level 强相关，所以用 noise level 作为 exploration capacity 的 proxy。

修改后的 policy loss 是：

\[
\mathcal{J}_{\text{policy}}(\theta)
=
\frac{1}{G}
\sum_{i=1}^{G}
\frac{1}{T}
\sum_{t=0}^{T-1}
\mathrm{Norm}(\sigma_t\sqrt{\Delta t})
\min
\left(
r_t^i(\theta)\hat{A}_t^i,
\mathrm{clip}(r_t^i(\theta),1-\epsilon,1+\epsilon)\hat{A}_t^i
\right)
\]

这里的 \(\mathrm{Norm}(\sigma_t\sqrt{\Delta t})\) 会被归一化，使所有 timestep 权重均值为 1。这样总 loss scale 不变，但不同时间点的相对优化强度改变。

\begin{center}
  \includegraphics[width=0.46\linewidth]{tempflowgrpo/fig6_noise_reward_correlation.png}
  \hfill
  \includegraphics[width=0.46\linewidth]{tempflowgrpo/fig6_scale_term.png}

  {\small \textbf{图 3：} 左：noise level 和 reward std 的相关性。右：标准 GRPO 的 scale term 与 noise-aware reweighting 后的 scale term 对比。图源：arXiv HTML。}
\end{center}

\subsubsection{3.1 理论推导在说明什么}

论文从 policy gradient 写起：

\[
\nabla_\theta \mathcal{J}(\theta)
=
\sum_{k=0}^{T-1}
\mathbb{E}
\left[
\nabla_\theta
\log p_\theta(x_{k-1}|x_k)
\hat{A}_k
\right]
\]

标准形式下，推导出的 timestep scale term 近似为：

\[
\sqrt{\frac{\Delta k(1-k)}{k}}
\]

加入 noise-aware reweighting 后，scale term 变成：

\[
\Delta k
\]

作者的解释是：标准 GRPO 的梯度贡献和真实 exploration capacity 不匹配，会让低噪声 refinement steps 在优化里占不合理的比重；noise-aware weighting 后，各 timestep 的梯度贡献更接近均衡。

\begin{summarybox}
\textbf{这里的重点不是“早期一定更大越好”。} 更准确地说，TempFlow-GRPO 想让优化强度匹配每个 timestep 的探索空间。早期有结构性探索，不能被低估；后期是细节修复，不能被过度扰动。
\end{summarybox}

\subsection{4. 方法三：Seed Group}

Trajectory branching 要比较不同 branches 的 reward。如果这些 branches 的初始噪声也不同，那 reward 差异会混入 initial noise 的影响，不能干净地归因到 branching point。

所以作者提出 seed group：

\[
\text{same prompt}
\quad + \quad
\text{same initial noise}
\quad + \quad
\text{different branch noise at } k
\]

这样 group 内 reward variance 更接近来自指定 timestep 的 exploration，而不是来自不同 \(x_T\) 的初始化差异。

实验配置里总 group size 保持为 24，默认是：

\[
4\ \text{initial seeds}
\times
6\ \text{branches}
=
24
\]

\begin{summarybox}
\textbf{这点对我们很有启发：} 如果要研究 credit assignment，group 的构造本身就是方法的一部分。不是随便采 \(G\) 条 trajectory，而是要控制住 prompt、initial noise、branching point，才能让 reward 差异有可解释性。
\end{summarybox}

\subsection{5. 实验结果}

\subsubsection{5.1 GenEval}

在 GenEval 上，论文报告：

\begin{center}
\begin{tabular}{lccc}
\toprule
方法 & step & overall & 备注 \\
\midrule
SD3.5-M base & - & 0.63 & base flow model \\
Flow-GRPO & 3800 & 0.88 & 同步数对比 \\
Flow-GRPO & 5600 & 0.95 & 更长训练 \\
TempFlow-GRPO & 3800 & 0.97 & 最优 \\
\bottomrule
\end{tabular}
\end{center}

论文还说 TempFlow-GRPO 约 2000 steps 就达到 0.95，而 Flow-GRPO 需要约 5600 steps。

\subsubsection{5.2 PickScore 和 FLUX}

在 PickScore benchmark 上：

\begin{itemize}[leftmargin=1.5em]
  \item 相比原始 Flow-GRPO，TempFlow-GRPO 大约提升 1.7\%。
  \item 相比改进后的 Flow-GRPO Prompt baseline，大约提升 1.0\%。
  \item 只需 100--200 training steps 就能达到 Flow-GRPO 的水平，300--400 steps 达到 Flow-GRPO Prompt 的水平。
\end{itemize}

在 FLUX.1-dev + HPSv3 上：

\begin{itemize}[leftmargin=1.5em]
  \item 约 80 steps 达到 Flow-GRPO 300 steps 的性能。
  \item Appendix 报告 300 steps 后相比 Flow-GRPO 约有 0.3 的性能提升。
  \item KL loss 更低、更稳定。
\end{itemize}

\subsubsection{5.3 和 DanceGRPO 的比较}

Appendix 里还和 DanceGRPO 做了比较。论文报告：

\begin{itemize}[leftmargin=1.5em]
  \item DanceGRPO 300 iterations 后 HPSv2 约 37.2。
  \item TempFlow-GRPO 280 iterations 达到 38.5。
  \item TempFlow-GRPO 约 150 iterations 就匹配 DanceGRPO 300 iterations 的最终效果，相当于约 \(2\times\) convergence speedup。
\end{itemize}

\subsubsection{5.4 Ablation}

\begin{center}
  \includegraphics[width=0.31\linewidth]{tempflowgrpo/fig9_branch_ablation.png}
  \hfill
  \includegraphics[width=0.31\linewidth]{tempflowgrpo/fig9_pickscore_ablation.png}
  \hfill
  \includegraphics[width=0.31\linewidth]{tempflowgrpo/fig9_geneval_ablation.png}

  {\small \textbf{图 4：} TempFlow-GRPO 的 ablation。左：不同 seed/branch 组合。中/右：trajectory branch、noise-aware weighting、seed group 的贡献。图源：arXiv HTML。}
\end{center}

几个关键结论：

\begin{itemize}[leftmargin=1.5em]
  \item 在总 group size 固定为 24 时，比较了 \(2\times 12\)、\(4\times 6\)、\(6\times 4\)。
  \item 更多 initial noises 有利于 early training 的快速提升。
  \item 更多 branches 在训练后期更好，因为同一个 branching point 下的探索比较更充分。
  \item 最终默认选择 \(4\times 6\)，平衡 initial diversity 和 branch comparison。
  \item 在 GenEval 上，noise-aware policy reweighting 从 0.82 提升到 0.92；trajectory branch 约带来 5\% gain；seed group 约带来 2\% gain。
\end{itemize}

\subsection{6. 成本问题}

TempFlow-GRPO 并不是免费午餐。由于每个 timestep 要做 \(K\)-branch exploration，采样成本高于 Flow-GRPO。论文在 Appendix A.6 里给了一个估计：

\[
K=10
\quad \Rightarrow \quad
\text{average branch factor}
\approx
\frac{9+\cdots+1}{10}
=
4.5
\]

也就是说，单个样本的采样过程可能比 Flow-GRPO 更贵。但论文认为它整体训练效率更高，因为 convergence 更快；在高分辨率任务上达到 Flow-GRPO 最终性能只需要约 33\%--50\% 的训练时间。

\begin{center}
  \includegraphics[width=0.46\linewidth]{tempflowgrpo/fig14_training_time_geneval.png}
  \hfill
  \includegraphics[width=0.46\linewidth]{tempflowgrpo/fig14_training_time_pickscore.png}

  {\small \textbf{图 5：} 训练时间对比。TempFlow-GRPO 单次采样更贵，但由于收敛更快，达到同等性能的 wall-clock time 更低。图源：arXiv HTML。}
\end{center}

\begin{warnbox}
\textbf{我的看法：} 这里的成本叙事要谨慎。TempFlow-GRPO 的优势是 \key{sample efficiency / convergence speed}，不是单步计算便宜。对图像任务这个 tradeoff 可能成立；但如果迁移到视频模型，多 branch 采样、decode、reward evaluation 的成本会被放大，需要重新评估。
\end{warnbox}

\subsection{7. 和 DenseGRPO / Chunk-GRPO 的关系}

\begin{center}
{\small
\setlength{\tabcolsep}{3pt}
\begin{tabular}{p{0.20\linewidth}p{0.25\linewidth}p{0.22\linewidth}p{0.22\linewidth}}
\toprule
方法 & reward / advantage & 更新粒度 & 核心代价 \\
\midrule
DenseGRPO & 对每个 \(x_t\) ODE readout，计算 \(\Delta R_t\) & step-level & full ODE readout 接近 \(O(T^2)\) \\
Chunk-GRPO & 主要仍用 terminal \(A_i\) & chunk-level ratio & credit 仍不是真正 chunk-specific \\
TempFlow-GRPO & ODE-SDE-ODE branch 后用最终 reward 得到 \(\hat{A}_k\) & step-level with temporal weighting & 多 branch exploration 成本高 \\
\bottomrule
\end{tabular}
}
\end{center}

三篇论文其实给了三种不同的 credit assignment 路径：

\[
\text{DenseGRPO：}
\quad
\text{readout-based local improvement}
\]

\[
\text{Chunk-GRPO：}
\quad
\text{chunk-level optimization granularity}
\]

\[
\text{TempFlow-GRPO：}
\quad
\text{branching-based counterfactual advantage}
\]

\begin{summarybox}
\textbf{新的认识：} step-level reward 不一定只能来自 \(\Delta R_t=R_{t-1}-R_t\)。还可以通过 controlled branching 构造：固定其他路径，只在一个 timestep 注入探索，再用最终 reward 差异反推该 timestep 的 advantage。
\end{summarybox}

\subsection{8. 局限与疑问}

\begin{enumerate}[leftmargin=1.5em]
  \item \textbf{它不是 DenseGRPO 式 improvement。} TempFlow-GRPO 的 \(R_k^i\) 是 branch outcome reward，不是相邻状态差分 \(R_{k-1}-R_k\)。因此它回答的是“在 \(k\) 探索会导致什么最终结果”，不是“从 \(x_k\) 到 \(x_{k-1}\) 本身让 reward 提升了多少”。
  \item \textbf{credit localization 依赖 controlled setup。} 如果初始 noise、prompt、ODE remainder 没控制好，reward 差异就不再只来自 branch noise。
  \item \textbf{多 branch 成本可能很高。} 对视频生成尤其明显，因为每个 branch 都可能意味着完整后续 ODE、VAE decode 和 reward scoring。
  \item \textbf{noise-aware weighting 是 proxy。} noise level 和 reward variance 相关，但不等于真实 semantic importance。某些任务里 late-stage detail 或 identity preservation 也可能很关键。
  \item \textbf{仍然依赖 outcome reward model。} 如果最终 reward model 本身偏了，branch advantage 也会跟着偏。
\end{enumerate}

\subsection{9. 对我们后续思考的启发}

这篇读完之后，我们可以把当前问题改得更深一点：

\[
\text{credit assignment 的单位是什么？}
\]

现在不只是 step vs chunk，还多了一个维度：

\[
\text{credit signal 是 readout difference，还是 controlled branching difference？}
\]

也就是说，未来方法不一定只是：

\[
\text{Chunk-GRPO} + \text{DenseGRPO}
\]

更有意思的方向可能是：

\[
\text{temporal segment}
\quad + \quad
\text{controlled branching}
\quad + \quad
\text{low-cost readout}
\]

例如，对一个 segment \([\tau_j,\tau_{j-1}]\)，不是机械做 DenseGRPO 式 boundary readout，也不是简单用 terminal reward，而是做 segment-level intervention：

\[
x_{\tau_j}
\xrightarrow{\mathrm{SDE\ segment}}
x_{\tau_{j-1}}
\xrightarrow{\mathrm{ODE}}
\hat{x}_0
\xrightarrow{R}
R_{\mathrm{seg}_j}
\]

然后同一个 segment 下比较多个 branches：

\[
\hat{A}_{\mathrm{seg}_j}^i
=
\frac{
R_{\mathrm{seg}_j}^i-\mathrm{mean}(\{R_{\mathrm{seg}_j}^g\})
}{
\mathrm{std}(\{R_{\mathrm{seg}_j}^g\})+\epsilon
}
\]

这会比“Chunk-GRPO + DenseGRPO”更自然，因为它把新问题定义成：

\begin{quote}
如何在合适的 temporal resolution 上构造可归因的 counterfactual reward signal？
\end{quote}

\begin{summarybox}
\textbf{读完要记住的三句话：}

1. TempFlow-GRPO 证明 step-level credit 不一定要靠中间 reward model，可以靠 controlled branching + final reward 构造。\\
2. 它把探索和评价分开：指定 timestep 用 SDE 探索，其余路径用 ODE 固定，再用 outcome reward 评价。\\
3. 它进一步说明 temporal credit assignment 的关键不是“越细越好”，而是要同时考虑 \key{归因可解释性、探索能力、更新强度和计算成本}。
\end{summarybox}
